\documentclass[11pt]{letter}

\usepackage[margin=0.5in]{geometry}
\usepackage{hyperref}

\begin{document}

\begin{letter}{}

\opening{Dear Hiring Team,}

I am writing to apply for the Software Application Developer, Generative AI position at BMO in Calgary. With a Master of Science in Computer Science from the University of Calgary and hands-on experience building backend APIs, Python/C++ data systems, SQL-backed services, and LLM-enabled automation workflows, I am excited by the opportunity to help integrate practical GenAI capabilities into enterprise applications.

What interests me about this role is the balance between software engineering and applied AI. The posting describes work with LLM services, retrieval patterns, AI-enabled workflows, APIs, cloud services, SQL, testing, deployment, and responsible validation of AI outputs. That combination fits the work I enjoy most: building useful systems around complex data and AI behavior, then making those systems reliable enough for real users and business decisions.

My graduate research with BigGeo, Mitacs, and the University of Calgary focused on building scalable Python and C++ software for large Earth observation datasets. I developed pipelines for data ingestion, preprocessing, storage, retrieval, and evaluation, and worked carefully with performance, reproducibility, and documentation. I improved one C++ processing workflow from 233 seconds to 26 seconds for a 10-tensor workload and reduced storage by 38\% on a scaled Canada-wide dataset. This experience strengthened the engineering discipline I would bring to enterprise AI systems: correctness, maintainability, and careful validation.

I also bring backend application experience. As a back-end developer at Asr Gooyesh Pardaz, I built Django REST APIs for NLP, text-to-speech, and speech-to-text services backed by PostgreSQL, implemented authentication workflows, integrated payment-service logic, and contributed to a chatbot system using GPT-style responses and FAQ data. That work gave me practical experience with APIs, relational databases, user-facing backend behavior, and troubleshooting service workflows.

More recently, I have been building an agentic job-search automation platform using LLM workflows, n8n, OpenClaw, Docker, Google Sheets, Google Drive, SearXNG, and Telegram API integrations. The project involves prompt-driven analysis, structured outputs, workflow automation, and practical decision support. While this is not a banking product, it is directly aligned with the kind of applied GenAI work described in this role: connecting LLM behavior to software workflows, validating outputs, and designing systems that are useful beyond a demo.

I have not worked professionally with Java, Spring Boot, Oracle, or production AWS services at the same depth as my Python, SQL, and backend experience, but I have a strong software foundation and have learned new technical environments quickly throughout my research and industry work. I am based in Calgary and would be glad to work in BMO's hybrid model while contributing to responsible, maintainable GenAI application development.

Thank you for considering my application. I would welcome the opportunity to discuss how my backend, AI workflow, SQL, data systems, and research software experience can contribute to BMO's Generative AI application work.

\closing{Sincerely,

Amir Mirzai Golpayegani

\href{mailto:amirgolpaa24@gmail.com}{amirgolpaa24@gmail.com}}

\end{letter}

\end{document}
